%% shortname: The Baseline Real Business Cycle Model
\documentclass[11pt]{article}
\usepackage{macrostyle}

\begin{document}

\noteheader{Business Cycle Theory}{The baseline Real Business Cycle model}{}


\section{Setup}

\runin{Notation.} Upper-case letters denote economy-wide averages that are taken as
given by individual households, whereas lower-case letters denote individual-specific
variables, which households control.

\topic{The representative firm's problem (FP)}
\begin{align}
  &\max\ \bigl\{\,Y_t-w_tN_t-r_tK_t\,\bigr\} \tag{1}\label{b:fp}\\
  \text{s.t.}\quad Y_t &= z_tF(K_t,A_tN_t) \tag{2}\label{b:prod}\\
  A_t &= A_0e^{xt} \tag{3}\label{b:A}\\
  \ln z_t &= \rho\ln z_{t-1}+\varepsilon_t,\qquad
     \varepsilon_t\sim\text{iid } N(0,\sigma^2) \tag{4}\label{b:z}
\end{align}
The first-order conditions give factor prices equal to marginal products:
\begin{align}
  \therefore\quad w_t &= z_tA_tF_2(t) \tag{5}\label{b:w}\\
  r_t &= z_tF_1(t) \tag{6}\label{b:r}
\end{align}

\topic{The representative household's problem (HP)}
\[ \max_{c_t,\,i_t,\,\ell_t,\,n_t}\ \E_0\Bigl\{\sum_{t=0}^{\infty}\beta^t u(c_t,\ell_t)\Bigr\} \]
\begin{align}
  \text{s.t.}\quad c_t+i_t &= w_tn_t+r_tk_t \tag{7}\label{b:bc}\\
  i_t &= k_{t+1}-(1-\delta)k_t \tag{8}\label{b:lom}\\
  n_t+\ell_t &= 1 \tag{9}\label{b:time}
\end{align}


\section{Dynamic programming with a Lagrange multiplier}

Substituting \eqref{b:lom} and \eqref{b:time} into \eqref{b:bc}, the dynamic programming
formulation of this problem is given by Bellman's equation, with a multiplier
$\lambda_t$ on the budget constraint:
\begin{align*}
  V(k_t;z_t)=\max_{c_t,\,k_{t+1},\,n_t}\Bigl\{\,& u(c_t,1-n_t)+\beta\E_tV(k_{t+1};z_{t+1})\\
  &+\lambda_t\bigl[w_tn_t+(r_t+1-\delta)k_t-c_t-k_{t+1}\bigr]\Bigr\}.
\end{align*}
Here $k_t$ (individual capital) and $z_t$ (technology) are the states; $c_t$, $n_t$ and
$k_{t+1}$ are the controls. Write $u_1(t)$, $u_2(t)$ for the partial derivatives of
$u(c,\ell)$ evaluated at date-$t$ choices, and $V'(t)=\partial V(k_t;z_t)/\partial k_t$.

\topic{First-order conditions}
\begin{align}
  \foc{c_t}{u_1(t)=\lambda_t} \tag{10}\label{b:fc}\\
  \foc{n_t}{u_2(t)=w_t\lambda_t} \tag{11}\label{b:fn}\\
  \foc{k_{t+1}}{\beta\E_tV'(t+1)=\lambda_t} \tag{12}\label{b:fk}
\end{align}

\topic{Envelope theorem}
\begin{align*}
  \foc{k_t}{V'(t)=\lambda_t(r_t+1-\delta)}
\end{align*}
Update \eqref{b:fc} and the envelope condition one period:
\begin{align}
  \foc{c_{t+1}}{u_1(t+1)=\lambda_{t+1}} \notag\\
  \foc{k_{t+1}}{V'(t+1)=\lambda_{t+1}(r_{t+1}+1-\delta)} \tag{13}\label{b:env}
\end{align}

\topic{Eliminate the multipliers}

\runin{Labour--leisure choice.} From \eqref{b:fc} and \eqref{b:fn}:
\begin{equation}
  u_2(t)=w_t\,u_1(t). \tag{14}\label{b:intra}
\end{equation}
This equates the marginal cost of working to the marginal benefit:
\begin{flatlist}
  \item $u_2(t)$ = marginal cost of working (foregone leisure);
  \item $w_t\,u_1(t)$ = marginal benefit of working (the wage $\times$ marginal utility
        from $c$).
\end{flatlist}

\runin{Intertemporal choice.} From \eqref{b:fc}, \eqref{b:fk} and \eqref{b:env}:
\begin{equation}
  u_1(t)=\beta\E_t\bigl\{u_1(t+1)\bigl[r_{t+1}+1-\delta\bigr]\bigr\}. \tag{15}\label{b:euler}
\end{equation}
This is the \emph{Euler equation}: it equates the marginal utility of consuming current
income to the discounted marginal utility you would receive by investing current income
and augmenting future consumption.

Alternative interpretation:
\begin{flatlist}
  \item LHS = marginal cost of saving (or investing);
  \item RHS = marginal benefit of saving (or investing).
\end{flatlist}


\section{Competitive equilibrium}

A competitive equilibrium consists of an allocation $\{c_t,n_t,k_{t+1},i_t\}$ for the
typical household, an allocation $\{N_t,K_t,Y_t\}$ for the typical firm, and a set of
prices $\{w_t,r_t\}$, such that:
\begin{flatlist}
  \item[(i)] the household's allocation solves HP, given prices;
  \item[(ii)] the firm's allocation solves FP, given prices;
  \item[(iii)] all markets clear:
\end{flatlist}
\begin{align}
  n_t &= N_t \tag{16}\label{b:mcn}\\
  k_t &= K_t \tag{17}\label{b:mck}\\
  c_t+i_t &= Y_t \tag{18}\label{b:mcy}
\end{align}


\section{The effects of a technology shock}

Suppose there is a positive technology shock in the current period:
$\varepsilon_t>0\Rightarrow z_t>1$ (given $z_{t-1}=1$).

\runin{From \eqref{b:z}:} $z_{t+1}$:\quad
$\E_t\{\ln z_{t+1}\}=\rho\ln z_t\ \Rightarrow\ \E_t\{z_{t+1}\}\uparrow$.

\runin{From \eqref{b:r}:} $r_{t+1}$:\quad
$\E_t\{r_{t+1}\}\uparrow\ \Rightarrow\ \E_t\{r_{t+1}+1-\delta\}\uparrow$.

\runin{From \eqref{b:euler}:} $k_{t+1}$, $i_t$:\quad RHS of \eqref{b:euler}
$\uparrow\ \Rightarrow\ k_{t+1}$ and $i_t\uparrow$.

$c_t$: save more and consume less (\emph{intertemporal substitution effect}, ISE).

\runin{From \eqref{b:intra}:}
\begin{flatlist}
  \item[(i)] $c_t\downarrow$ via ISE $\Rightarrow$ RHS of \eqref{b:intra}$\uparrow$
        $\Rightarrow n_t\uparrow,\ \ell_t\downarrow$.\\
        Intuition: work more and save more.
  \item[(ii)] $w_t\uparrow$ via \eqref{b:w} $\Rightarrow$ RHS of \eqref{b:intra}$\uparrow$
        $\Rightarrow n_t\uparrow,\ \ell_t\downarrow$.\\
        Intuition: substitution effect $\to$ work more when $w$ is high.
\end{flatlist}

\runin{From \eqref{b:prod}:} $Y_t\uparrow$ because $n_t\uparrow$ and $z_t\uparrow$.

Since $Y_t\uparrow$ (and $w_t$, $r_t$), $c_t$ may $\uparrow$. This is called the
\emph{income effect} (IE).

\topic{Summary}
\[ \begin{aligned}
   \text{ISE}&\Rightarrow c_t\downarrow\\
   \text{IE}&\Rightarrow c_t\uparrow
   \end{aligned}
   \qquad
   \left\{\ \parbox{0.5\linewidth}{Intuition: you save a higher fraction of $Y_t$
   ($\therefore c_t\downarrow$) but you have more $Y_t$ ($\therefore c_t\uparrow$).}\right. \]

\runin{Note.} The IE also affects $n_t$: if $c_t\uparrow$ via IE then RHS of
\eqref{b:intra}$\downarrow\ \Rightarrow n_t\downarrow$ and $\ell_t\uparrow$.
Intuition: leisure is a normal good.

\topic{Results}
\[ z_t\uparrow\ \Rightarrow\ Y_t\uparrow,\ i_t\uparrow,\
   c_t\uparrow\ (\text{if IE}>\text{ISE}),\ w_t\uparrow,\
   r_t+1-\delta\uparrow,\ n_t\uparrow\ (\text{if SE}+\text{ISE}>\text{IE}). \]
$\hookrightarrow$ Everything could be procyclical, just like the data. But can the model
match the \emph{size} of the observed correlations? Can it match the observed
\emph{volatility}? We need to solve the model and run a simulation.


\section{Solving the model}

Before we can run a simulation we need to solve the model and set parameter values.
Solving the model involves finding the \emph{policy functions} (or \emph{decision rules})
that map the current-period state variable and the current state of technology into the
control variables and the future value of the state variable, i.e.
\begin{align*}
  c_t &= C(k_t,z_t)\\
  k_{t+1} &= K(k_t,z_t)\\
  &\ \ \vdots\\
  &\ \ \text{etc.}
\end{align*}

\runin{Problem.} As is generally the case with non-linear dynamic systems, the model has
(in general) no closed-form analytical solution. Like the Ramsey model, we will solve the
general version of the RBC model by performing a \emph{log-linear approximation around
the steady state}. Before we do this, let's consider a special case of the model that has
an analytical solution. Assume:
\begin{align*}
  u(c,\ell) &= \omega\ln c+(1-\omega)\ln\ell\\
  Y &= zN^{1-\alpha}K^{\alpha}\\
  \delta &= 100\%
\end{align*}


\begin{transcriptnotes}
  \item Source: handwritten notes \texttt{RBC1.pdf} (6 pages). Equation numbers
        (1)--(18) follow the original; the notes stop after stating the assumptions of
        the special case (continued next lecture).
  \item The Bellman equation is written in the Lagrange-multiplier format, with the
        multiplier term entered as $+\lambda_t[\text{resources}-\text{uses}]$. The lecture
        wrote the equivalent $-\lambda_t[c_t+k_{t+1}-w_tn_t-(r_t+1-\delta)k_t]$; all FOCs
        are unchanged. The ``update'' and ``eliminate the multipliers'' steps are made
        explicit.
  \item (11): the original reads $u_2(t)=w_{\scriptscriptstyle 2}\lambda_t$; corrected
        to $w_t\lambda_t$.
  \item Policy functions: the original says they ``map the control variables and the
        future value of the state variable into the current-period state variable and
        technology''; reworded in the direction of the map, $(k_t,z_t)\mapsto(c_t,k_{t+1})$,
        consistent with the formulas $c_t=C(k_t,z_t)$, $k_{t+1}=K(k_t,z_t)$.
  \item ``SE'' in the Results line is the (intratemporal) wage substitution effect of
        item (ii) under \eqref{b:intra}.
\end{transcriptnotes}

\end{document}
